\documentclass[UTF8,a4paper,11pt]{ctexart}
\usepackage[a4paper,top=22mm,bottom=20mm,left=20mm,right=20mm,headheight=15pt]{geometry}
\usepackage{fontspec}
\usepackage{graphicx}
\usepackage{booktabs}
\usepackage{tabularx}
\usepackage{array}
\usepackage[table]{xcolor}
\usepackage{enumitem}
\usepackage{fancyhdr}
\usepackage{hyperref}
\setmainfont{Noto Serif CJK SC}
\setsansfont{Noto Sans CJK SC}
\setmonofont{DejaVu Sans Mono}
\setCJKmainfont{Noto Serif CJK SC}
\setCJKsansfont{Noto Sans CJK SC}
\definecolor{brand}{HTML}{185B46}
\definecolor{branddark}{HTML}{173F33}
\definecolor{brandpale}{HTML}{E5F0EB}
\definecolor{brandgray}{HTML}{566B62}
\hypersetup{unicode=true,colorlinks=true,linkcolor=brand,urlcolor=brand,pdftitle={ESG ClaimGuard：面向可持续披露的证据约束式智能预审系统}}
\pagestyle{fancy}
\fancyhf{}
\fancyfoot[C]{\thepage}
\renewcommand{\headrulewidth}{0pt}
\setlength{\parindent}{2em}
\setlength{\parskip}{0.35em}
\renewcommand{\arraystretch}{1.35}
\newcolumntype{Y}{>{\raggedright\arraybackslash}X}
\setlist{leftmargin=2.2em,itemsep=0.25em,topsep=0.35em}
\ctexset{
  section={format=\Large\bfseries\color{branddark},beforeskip=1.4em,afterskip=0.7em},
  subsection={format=\large\bfseries\color{brand},beforeskip=1.1em,afterskip=0.5em},
  subsubsection={format=\normalsize\bfseries\color{brand},beforeskip=0.9em,afterskip=0.4em}
}
\emergencystretch=3em
\raggedbottom
\begin{document}
\begin{center}
{\fontsize{23}{30}\selectfont\bfseries\sffamily\color{branddark} ESG ClaimGuard：面向可持续披露的证据约束式智能预审系统\par}

\vspace{10mm}
\end{center}
\subsection{摘要}

针对 ESG 报告版面复杂、指标表达分散、生成式抽取难以审计的问题，本文提出 ESG ClaimGuard。方法以文档版面解析、指标候选召回、受约束大模型抽取和确定性证据门组成分层流水线，并将定量结果区分为直接读取、表达式推导和单位归一。系统在 200 份报告、10,528 页和 65 项指标上生成 13,000 条唯一结果，其中 found 7,688、missing 5,312、error 0。所有 found 均通过原始解析区块精确子串校验，3,214 条定量结果均具有值与单位来源。本文强调这些结果证明工程完整性和证据可追溯性，不等同于独立准确率；本作品不把未经独立人工评测的 Precision、Recall 或 F1 作为成果声明。\par

\textbf{关键词：} ESG；信息抽取；检索增强生成；证据追溯；人机协同\par

\subsection{问题定义}

给定报告集合 R 与指标集合 I，系统为每个二元组 (r, i) 输出 found 或 missing。found 必须携带页码、区块、原始引文及结构化值；missing 表示在当前解析和召回范围内没有满足证据合同的候选，不等价于事实上的未披露。证据门的纯文本约束为：evidence\_quote q 必须是 canonical block b(r, p) 的连续精确子串，其中 p 是报告 r 的页码。\par

\subsection{方法}

流水线见图 1。MinerU 建立页级版面区块；候选召回限制模型上下文；Qwen3.6 输出受约束 JSON；确定性验证检查唯一键、状态、引文、页码、区块和定量来源。若数值为直接读取，记录 value\_origin 和 unit\_origin；若为推导值，同时保存 derivation\_expression 与 derivation\_inputs。任何强制字段失败均保守退化为 missing 或进入人工复核。\par

\begin{center}

\includegraphics[width=0.96\linewidth,height=0.72\textheight,keepaspectratio]{\detokenize{latex/figures/fig_pipeline.png}}

\par\small\color{brandgray} 图 1　ESG ClaimGuard 顺序流水线

\end{center}

\subsection{数据与实验协议}

冻结队列由 \texttt{input\_manifest.csv} 选取 200 份公开 ESG 类报告，共 10,528 页；项目目录中的其他 PDF 不属于该正式队列。指标池含环境 25、社会 20、治理 20 项，形成 13,000 个 report×indicator 任务。正式运行使用 MinerU2.5-Pro-2605-1.2B 和 Qwen3.6-27B Q4\_K\_M。本项目将 Qwen 官方仓库中的 Qwen3.6-27B 以 Q4\_K\_M 量化形态部署，仅使用文本输入和 8,192 token 运行上下文；模型身份与许可依据官方模型卡记录 [2]。209 条被重试熔断器保留的记录由 Codex Agent 模拟人工工程复核，其中 found 193、missing 16；该集合不作为独立测试集。\par

本文的验收范围是覆盖规模、证据合同、运行稳定性和时延。209 条工程复核记录仅用于修复与一致性检查，不作为独立人工测试集，因此不据此计算或宣称 Precision、Recall、F1。\par

\subsection{结果}

正式结果为 found 7,688、missing 5,312、error 0，唯一键 13,000/13,000。E/S/G 的 found 数分别为 2,363、2,550、2,775；报告级 found 数最小 0、中位数 42、最大 57。运行级生成调用为 10,015 次，总墙钟 27,928.801 秒；最终结果中 11,280 行保留正 \texttt{elapsed\_seconds}，该行级样本的中位数为 2.765 秒、P95 为 3.683 秒。\par

图 2 将运行级生成调用数与有 \texttt{elapsed\_seconds} 的结果行明确分开：直方图中的 11,280 个样本是有耗时记录的结果行，不是 10,015 次生成调用。两者来自不同统计层级，不能互换。\par

\begin{center}

\includegraphics[width=0.96\linewidth,height=0.72\textheight,keepaspectratio]{\detokenize{latex/figures/fig_inference.png}}

\par\small\color{brandgray} 图 2　运行级调用卡片与有耗时记录的结果行分布

\end{center}

证据合同方面，7,688 条 found 全部满足原始精确子串与页码区块追溯，失败 0。3,214 条定量结果的 value origin 为直接读取 3,089 条、明确推导 125 条；unit origin 为原文单位 3,105 条、单位归一或推断 109 条。这些计数由文档渲染器从正式 \texttt{validation.json} 注入，不在论文中固化旧值。上述“全部通过”仅表示内部字段和追溯合同完整，不能作为语义准确率。\par

\begin{center}

\includegraphics[width=0.96\linewidth,height=0.72\textheight,keepaspectratio]{\detokenize{latex/figures/fig_evidence_gate.png}}

\par\small\color{brandgray} 图 3　证据完成门、value origin 与 unit origin 分面

\end{center}

\subsection{讨论}

系统的主要贡献是把生成式抽取限制在可复核证据链中。相比自由问答，评审者可以定位原文并复算推导值；相比关键词，系统能处理同义语义；相比普通 RAG，缺失与定量来源成为显式合同。局限包括 OCR 误差、跨页表格、隐含范围和行业口径差异，因此本作品将结论限定在工程完整性与证据可追溯性。\par

\subsection{结论}

ESG ClaimGuard 已完成文档解析、指标抽取、证据验证、人工处置和底稿交付的端到端实现。正式运行覆盖 200 份报告与 13,000 个 report×indicator 任务，验证结果表明系统能够稳定生成带页码、区块、原文和定量来源的结构化候选。\par

该系统的核心价值在于把生成式模型能力转化为可追溯、可复核、可处置的业务任务，并以确定性证据合同控制模型输出边界。其可信主张是工程完整性和证据可追溯性，而非未经验证的准确率领先，符合审计场景对证据和责任的要求。\par

综上，ESG ClaimGuard 已形成完整的可持续披露一致性预审系统，完成既定研究目标与工程交付。\par

\subsection{参考文献}

[1] Wang B. et al. MinerU: An Open-Source Solution for Precise Document Content Extraction. arXiv:2409.18839, 2024.\newline [2] Qwen Team. Qwen3.6-27B Model Card. Hugging Face model repository, 2026. https://huggingface.co/Qwen/Qwen3.6-27B.\newline [3] Yang A. et al. Qwen3 Technical Report. arXiv:2505.09388, 2025.\newline [4] Gao Y. et al. Retrieval-Augmented Generation for Large Language Models: A Survey. arXiv:2312.10997, 2023.\newline [5] Fan W. et al. A Survey on RAG Meeting LLMs. KDD, 2024.\par
\end{document}
